\documentclass[10pt,a4paper]{amsart}
\usepackage[margin=19mm]{geometry}
\usepackage{amsmath,amssymb,mathtools}
\usepackage[scr=rsfs,cal=euler]{mathalfa}
\usepackage{xfrac}
\usepackage[unicode,hidelinks,bookmarks=false]{hyperref}
\newcommand{\ie}{i.e.\ }
\newcommand{\cf}{cf.\ }
\newcommand{\resp}{respectively\ }
\newcommand{\Subsection}{\subsection}
\providecommand{\nobreakdash}{\nobreak\hbox{-}}
\setlength{\emergencystretch}{2em}
\multlinegap0pt
\usepackage{amssymb}
\usepackage[shortlabels]{enumitem}
\usepackage{xcolor}
\usepackage{mathtools}
\usepackage{bbm}
\usepackage{dsfont}
\usepackage{tikz}
\usepackage{booktabs}
\usepackage{natbib}
\newtheorem{lemma}{Lemma}
\usepackage{cleveref}
\crefname{lemma}{Lemma}{Lemmas}
\def\NN {\mathbb{N}}
\def\RR {\mathbb{R}}
\def\Re{{\rm Re}}
\newcommand{\cO}{\mathcal{O}}
\newcommand{\cQ}{\mathcal{Q}}
\newcommand{\cS}{\mathcal{S}}
\newcommand{\cT}{\mathcal{T}}
\newcommand{\cV}{\mathcal{V}}
\newcommand{\cZ}{\mathcal{Z}}
\def\dd{{\rm d}}
\def\de{{\partial}}
\def\e{{\rm e}} 
\def\eps{\varepsilon}
\def\epsilon{\varepsilon}
\newcommand{\mrM}{\mathrm{M}}
\newcommand{\mrm}{\mathrm{m}}
\title{Viscous vortex crystals}
\author{Michele Dolce, Martin Donati}
\date{Source: arXiv:2602.23134v1 (CC BY 4.0)}
\begin{document}
\maketitle

\noindent\emph{Source and licence.} Viscous vortex crystals. Version arXiv:2602.23134v1 (CC BY 4.0), distributed by arXiv under the Creative Commons Attribution 4.0 International licence, http://creativecommons.org/licenses/by/4.0/, as recorded in the arXiv licence metadata for this exact version and shown on the abstract page https://arxiv.org/abs/2602.23134v1. Full licence notice: \texttt{licenses/arxiv-2602.23134-CC-BY-4.0.txt}.\par\smallskip

\noindent\textit{Editorial scope.} Counted unit: the article's lemma lem:expVel (source0.tex lines 3146--3155). The article's complete original proof of it is delivered with it (source0.tex lines 3156--3196, one \texttt{proof} environment of 41 lines). No mathematics is written, repaired or re-derived: the statement and the proof are the article's own lines, in the article's own notation, and no retained line is substituted or re-flowed.  Retained uncounted as the source-local context the proof needs: lines 3039--3061.  Nothing was omitted from any retained span, so the package carries no omission record.  No retained line contains a citation, so the wrapper carries no bibliography and no bibliography entry.  No retained line contains a figure environment, an image-inclusion command or drawing code, so no image is delivered and the package declares no asset dependency.  Editorial formatting only, in addition: an AMS \texttt{amsart} preamble that restates the article's own declarations and macros used by the retained lines, namely the environments \texttt{lemma} on separate counters, together with the standard packages those lines need.  The retained environments print in this document as Lemma 1 in order of appearance, whereas the article prints the same items with the numbering of its own full document.  The complete native source of the article as distributed in the arXiv e-print for this exact version is bundled unchanged as \texttt{sources/195372/source0.tex}.
\begin{lemma}\label{lem:Expansions}
    Let $f \in \cZ$,  $\varphi = \Delta^{-1}f$, $0<z_{\min}\leq |Z|$ and $M\in \NN$. Then, the translation operator admits the expansion
    \begin{equation}
        \cT_{\eps,Z}\varphi(\xi) \,=\, \frac{{\rm M}(f)}{2\pi}\log\Big(\frac{|Z|}{\eps d}\Big)
    +\sum_{n=1}^M\sum_{k=0}^{n} \eps^nc_{n,k} \Re\Big(\frac{\xi^{n-k}}{Z^n}{\rm m}^k(f)\Big) + \cO_{\cS_*}\bigl(\epsilon^{M+1}\bigr)\,.
    \end{equation}
    The polygon-center operator admits the expansion
    \begin{equation}
  \begin{split}
      \cV_{\eps,Z}\varphi(\xi) \,=\,& \mrM(\Omega)\frac{N}{2\pi}\log\Big(\frac{|Z|}{\eps d}\Big)
  +N\sum_{n=1}^M (-1)^n\eps^nc_{n,n} \Re\Big(\frac{1}{Z^n}{\rm m}^n(f)\Big)\\
&+N\sum_{n=N}^M\sum_{\substack{\ell \in \mathbb{N}\setminus \{0\}\\
\ell N\leq M}}(-1)^n\eps^n c_{n,n-\ell N} \Re\Big(\frac{\xi^{\ell N}}{Z^n}\mrm^{n-\ell N}(f)\Big)+\cO_{\cS_*}\bigl(\epsilon^{M+1}\bigr)\,.
        \end{split}
\end{equation}
    The polygon-polygon operator admits the expansion
\begin{equation}
  \begin{split}
\cQ_{\eps,Z}\varphi(\xi) \,=&\, \frac{\mrM(\Omega)}{2\pi}\sum_{l=1}^{N-1}\log\Big(\frac{|Z_l|}{\eps d}\Big)
  +\sum_{n=1}^M \sum_{k=0}^n\epsilon^nc_{n,k}{\rm S}_{n,n-k}\Re\Bigg(\frac{\xi^{n-k}}{Z^n}\mrm^k(f)\Bigg)+\cO_{\cS_*}(\eps^{M+1}).
        \end{split}
\end{equation}
\end{lemma}

\begin{lemma}
\label{lem:expVel}
    Let $f,h\in \cZ$, $\varphi=\Delta^{-1}f$ and assume that $0<z_{\min}\leq |Z|$. Then, for all $M\in \mathbb{N}$ one has the expansions
    \begin{align}
\label{eq:expspeed}&\int_{\RR^2}\nabla^\perp(\cT_{\eps,Z}\varphi)h\,\dd \xi=\sum_{n=1}^M\eps^n\frac{Z^\perp}{|Z|^{2}}\frac{1}{\overline{Z^{n-1}}}\sum_{k=0}^{n-1}(n-k)c_{n,k}\overline{\mrm^k(f)}\overline{\mrm^{n-k-1}(h)}+\cO(\eps^{M+1})    \end{align}
and
\begin{equation}
    \int_{\RR^2}\nabla^\perp(\cQ_{\eps,Z}\varphi)h\,\dd \xi=\sum_{n=1}^M\eps^n\frac{Z^\perp}{|Z|^{2}}\frac{1}{\overline{Z^{n-1}}}\sum_{k=0}^{n-1}(n-k)c_{n,k}{\rm S}_{n,n-k}\overline{\mrm^k(f)}\overline{\mrm^{n-k-1}(h)}+\cO(\eps^{M+1}) .
\end{equation}
\end{lemma}

\begin{proof}
  Appealing to \cref{lem:Expansions}, we know that 
    \begin{align}
\nabla^\perp\cT_{\eps,Z}\varphi=\sum_{n=1}^M\sum_{k=0}^n\eps^nc_{n,k}\nabla^\perp\Re\bigg(\frac{\xi^{n-k}}{Z^n}\mrm^k(f)\bigg)+\cO_{\cS_*}(\eps^{M+1}).
    \end{align}
    To compute the orthogonal gradient, we pass to polar coordinates $\xi=r\e^{i\theta}$ (keeping the complex number notation) and we note that
    \begin{equation}
        \nabla^\perp =i(\de_1+i\de_2)\to i\e^{i\theta}(\de_r+i\frac{1}{r}\de_\theta).
    \end{equation}
    Then, denoting 
    \begin{align}
        Z=|Z|\e^{i\phi_Z}, \qquad \mrm^k(f)=|\mrm^k(f)|\e^{i\phi_k},
    \end{align}
we get    \begin{align}
    \Re\bigg(\frac{\xi^{n-k}}{Z^n}\mrm^k(f)\bigg)\to \frac{|\mrm^k(f)|}{2|Z|^n}r^{n-k}\big(\e^{i(n-k)\theta +i(\phi_k-n\phi_Z)}+\e^{-i(n-k)\theta -i(\phi_k-n\phi_Z)}\big).
    \end{align}
    Thus, applying the $\nabla^\perp$ we see that the right-hand side becomes 
    \begin{align}
    &i(n-k)\frac{|\mrm^k(f)|}{2|Z|^n}r^{n-k-1}\big(\e^{i(n-k+1)\theta +i(\phi_k-n\phi_Z)}+\e^{-i(n-k-1)\theta -i(\phi_k-n\phi_Z)}\big)\\
    &-i(n-k)\frac{|\mrm^k(f)|}{2|Z|^n}r^{n-k-1}\big(\e^{i(n-k+1)\theta +i(\phi_k-n\phi_Z)}-\e^{-i(n-k-1)\theta -i(\phi_k-n\phi_Z)}\big)\\
    &=i(n-k)\frac{|\mrm^k(f)|}{|Z|^n}r^{n-k-1}\e^{-i(n-k-1)\theta -i(\phi_k-n\phi_Z)}.
    \end{align}
    Going back to Cartesian coordinates, note that
    \begin{align}
        r^{n-k-1}\e^{-i(n-k-1)\theta}\to \overline{\xi^{n-k-1}}, \qquad |\mrm^k(f)|\e^{-i\phi_k}\to \overline{\mrm^{k}(f)}
    \end{align}
    and
    \begin{align}
        i\frac{1}{|Z|^n}\e^{in\phi_Z}=i\frac{|Z|}{|Z|^{n+1}}\e^{in\phi_Z}\to \frac{Z^\perp}{|Z|^{2}}\frac{1}{\overline{Z^{n-1}}}.
    \end{align}
    Therefore, we have computed that 
    \begin{align}
\nabla^\perp\Re\bigg(\frac{\xi^{n-k}}{Z^n}\mrm^k(f)\bigg)=(n-k)\frac{Z^\perp}{|Z|^{2}}\frac{1}{\overline{Z^{n-1}}}\overline{\mrm^{k}(f)}\overline{\xi^{n-k-1}}.
    \end{align}
    Since $h\in \cZ$, we then easily deduce that 
    \begin{align}
        \int_{\RR^2}(\nabla^\perp\cT_{\eps,Z}\varphi)h\, \dd\xi=\sum_{n=1}^M\eps^n\frac{Z^\perp}{|Z|^{2}}\frac{1}{\overline{Z^{n-1}}}\sum_{k=0}^{n-1}(n-k)c_{n,k}\overline{\mrm^k(f)}\overline{\mrm^{n-k-1}(h)}+\cO(\eps^{M+1}),
    \end{align}
    which proves \cref{eq:expspeed}.
    The proof for $\cQ_{\eps,Z}$ is now straightforward using the same arguments.
\end{proof}

\end{document}
